\chapter{Viral infections}
\index{Viral infections}

\section*{Definition and Overview}
A virus is a submicroscopic infectious agent that can only replicate inside the living cells of an organism. Structurally, viruses are remarkably simple yet highly efficient biological entities. Unlike bacteria, fungi, or parasites, they lack the cellular machinery required for independent metabolism, protein synthesis, and self-reproduction. Consequently, they are classified as obligate intracellular parasites. A typical virion (a complete, infective viral particle) consists of a core of genetic material, which can be either deoxyribonucleic acid (DNA) or ribonucleic acid (RNA), surrounded by a protective protein shell called a capsid. Some viruses also possess an outer lipid membrane, known as an envelope, which is derived from the host cell's membrane during the budding process.

The scope of viral infections is vast, spanning across all domains of life. In humans, viral infections are responsible for a broad spectrum of clinical manifestations, ranging from mild, self-limiting illnesses such as the common cold (primarily caused by rhinoviruses) to severe, life-threatening systemic diseases like Ebola virus disease, rabies, and Human Immunodeficiency Virus (HIV/AIDS). Viral pathogens can target virtually every organ system, including the respiratory tract (e.g., influenza, respiratory syncytial virus, coronaviruses), the gastrointestinal tract (e.g., rotavirus, norovirus), the central nervous system (e.g., enteroviruses, herpes simplex virus), and the liver (e.g., hepatitis A, B, C, D, and E viruses).

The significance of studying and managing viral infections lies not only in their immediate clinical impact but also in their epidemiological volatility. Viruses are prone to genetic variations through mechanisms such as point mutations (antigenic drift) and genomic reassortment (antigenic shift). These evolutionary processes can lead to the emergence of novel viral strains with altered transmissibility, pathogenicity, or resistance to existing medical interventions. The rapid global spread of viruses, exemplified by historical influenza pandemics and the SARS-CoV-2 pandemic, highlights the critical intersection of virology, clinical medicine, and global public health. This chapter aims to provide a comprehensive, authoritative reference for understanding the biological, clinical, and preventative dimensions of human viral infections.

\section*{Epidemiology}
The epidemiology of viral infections is characterized by its high prevalence, seasonal variability, and significant burden on global healthcare systems. Viral diseases represent a major cause of morbidity and mortality worldwide, affecting billions of individuals annually. Acute respiratory infections of viral etiology are particularly ubiquitous. For instance, seasonal influenza is estimated to cause 3 to 5 million cases of severe illness and up to 650,000 respiratory deaths globally each year. Similarly, viral gastroenteritis represents a leading cause of childhood mortality, particularly in low- and middle-income countries where access to sanitation and clean water is limited.

Affected populations and demographics vary significantly based on the specific viral pathogen, transmission dynamics, and host vulnerability. While individuals of all ages can contract viral infections, severe manifestations and complications are disproportionately observed in specific vulnerable groups. These include:
\begin{itemize}
    \item Infants and young children, whose immune systems are immunologically naive and still developing.
    \item Older adults (typically aged 65 years and older), who experience immunosenescence—a gradual decline in immune system function associated with aging.
    \item Pregnant women, who undergo physiological and immunological adaptations that increase susceptibility to certain infections, such as influenza and hepatitis E.
    \item Immunocompromised individuals, such as patients undergoing chemotherapy, organ transplant recipients on immunosuppressive therapy, and individuals with advanced HIV infection.
\end{itemize}

The geographical distribution of viral infections is heavily influenced by climate, vector availability, and public health infrastructure. Vector-borne viral infections, such as dengue, Zika, chikungunya, and yellow fever, are primarily endemic to tropical and subtropical regions where their mosquito vectors (such as \textit{Aedes aegypti}) thrive. Conversely, respiratory viruses, such as influenza and respiratory syncytial virus (RSV), exhibit pronounced seasonal patterns in temperate regions, peaking during the winter months. This seasonality is driven by factors such as low humidity, colder temperatures (which enhance viral stability in the air), and increased indoor crowding.

\section*{Aetiology and Pathophysiology}
The initiation and progression of viral infections involve a complex interplay between viral factors and the host's cellular defense mechanisms. Understanding the molecular pathways of viral pathogenesis is essential for developing targeted therapeutic and preventive strategies.
\begin{itemize}
    \item \textbf{Viral Classification and Diversity}: Viruses are broadly classified based on their nucleic acid composition (double-stranded or single-stranded DNA or RNA), capsid symmetry (icosahedral, helical, or complex), and the presence or absence of a lipid envelope. RNA viruses generally exhibit higher mutation rates than DNA viruses due to the lack of proofreading activity in viral RNA polymerases, facilitating rapid adaptation and immune evasion.
    \item \textbf{Transmission Pathways}: Viral pathogens enter the human host through distinct portals of entry. Common pathways include respiratory droplets and aerosols (e.g., influenza, SARS-CoV-2), the fecal-oral route (e.g., enteroviruses, norovirus), direct mucosal or sexual contact (e.g., Herpes Simplex Virus, Human Immunodeficiency Virus), blood-borne inoculation (e.g., Hepatitis B and C), and vector-borne transmission via arthropods (e.g., West Nile virus).
    \item \textbf{Cellular Attachment and Entry}: The infectious process begins with the specific binding of viral surface proteins (ligands) to complementary receptors on the host cell membrane. This interaction determines viral tissue tropism—the specific cells and tissues a virus can infect. For example, the gp120 glycoprotein of HIV binds to the CD4 receptor on T-helper lymphocytes, while the spike protein of SARS-CoV-2 binds to the angiotensin-converting enzyme 2 (ACE2) receptor on respiratory epithelial cells. Following attachment, the virus enters the cell via endocytosis or direct membrane fusion.
    \item \textbf{Intracellular Replication and Assembly}: Once inside the host cell, the viral capsid undergoes uncoating, releasing the viral genome. The virus co-opts the host's metabolic machinery, enzymes, and ribosomes to transcribe viral messenger RNA, translate viral proteins, and replicate the viral genome. The newly synthesized components are assembled into progeny virions, which are then released from the cell through lysis (rupturing the cell) or budding (acquiring an envelope from the host cell membrane).
    \item \textbf{Pathogenesis of Host Injury}: Viral-induced cellular damage occurs through direct and indirect pathways. Direct cytopathic effects (CPE) include cellular lysis, disruption of cellular protein synthesis, and the induction of apoptosis. Indirect damage is mediated by the host's immune system; an overactive immune response can lead to extensive tissue damage (immunopathology). A prime example is the ``cytokine storm,'' characterized by the uncontrolled release of pro-inflammatory cytokines, leading to systemic inflammatory response syndrome and multi-organ failure.
    \item \textbf{Immune Evasion Mechanisms}: To survive and replicate, many viruses have evolved strategies to bypass host defenses. These include downregulating Major Histocompatibility Complex (MHC) class I molecules to evade cytotoxic T-lymphocyte recognition, expressing viral proteins that inhibit interferon signaling pathways, and establishing latency, where the viral genome remains dormant within the host cell (e.g., sensory neurons for Herpes Simplex Virus) without active viral replication.
\end{itemize}

\section*{Clinical Features}
The clinical presentation of viral infections is highly variable, ranging from asymptomatic colonization to acute, life-threatening systemic syndromes. The presentation depends on the tissue tropism of the virus, the viral load, and the robustness of the host's immune response.
\begin{itemize}
    \item \textbf{Incubation and Prodromal Phases}: The incubation period is the time elapsed between viral exposure and the onset of clinical symptoms, which can range from days (e.g., rhinovirus) to years (e.g., HIV). The prodromal phase follows, characterized by non-specific, constitutional symptoms as the innate immune system initiates a generalized inflammatory response. These symptoms typically include mild fever, malaise, fatigue, generalized myalgia, and headache.
    \item \textbf{Acute Systemic Manifestations}: As the virus replicates and disseminates, systemic symptoms become more pronounced. High-grade fever, rigors, arthralgia, and profound fatigue are common. These symptoms are largely driven by circulating cytokines, such as interferons and interleukins, which reset the hypothalamic thermostat and induce muscle and joint discomfort.
    \item \textbf{Respiratory Signatures}: Infections targeting the respiratory tract present with symptoms categorized by anatomical location. Upper respiratory tract signs include rhinorrhea, sneezing, nasal congestion, pharyngitis, and hoarseness. Lower respiratory tract involvement presents with cough (productive or non-productive), sputum production, tachypnoea, wheezing, and dyspnoea. Auscultation may reveal crackles or decreased breath sounds.
    \item \textbf{Gastrointestinal Presentations}: Viral gastroenteritis classically presents with the acute onset of watery, non-bloody diarrhoea, accompanied by nausea, vomiting, abdominal cramping, and low-grade fever. Dehydration is the primary clinical concern, marked by dry mucous membranes, decreased skin turgor, tachycardia, and oliguria.
    \item \textbf{Dermatological and Mucocutaneous Findings}: Many viral infections present with distinctive exanthems (skin rashes) or enanthems (mucous membrane lesions). These include maculopapular rashes (e.g., measles, rubella), vesicular rashes (e.g., varicella-zoster virus, herpes simplex virus), and petechial or purpuric rashes (associated with viral hemorrhagic fevers).
    \item \textbf{Neurological Deficits}: Neurotropic viruses can cause inflammation of the meninges (meningitis) or brain parenchyma (encephalitis). Clinical features include severe headache, photophobia, phonophobia, nuchal rigidity (neck stiffness), altered mental status, confusion, seizures, and focal neurological deficits.
\end{itemize}

\section*{Differential Diagnosis}
Given that the clinical features of viral infections frequently overlap with those of other infectious and non-infectious conditions, a rigorous differential diagnosis is essential to ensure appropriate management.
\begin{itemize}
    \item \textbf{Bacterial Infections}: Bacterial pathogens often present with symptoms similar to viral illnesses, but distinguishing between them is critical to prevent the inappropriate use of antibiotics. For example, bacterial pharyngitis (e.g., \textit{Streptococcus pyogenes}) must be differentiated from viral pharyngitis, and bacterial meningitis must be rapidly distinguished from viral (aseptic) meningitis, as the former requires immediate intravenous antibiotic therapy to prevent mortality.
    \item \textbf{Fungal and Parasitic Infections}: In patients presenting with persistent fevers, atypical pneumonia, or chronic gastrointestinal symptoms, fungal pathogens (e.g., \textit{Pneumocystis jirovecii}, \textit{Histoplasma capsulatum}) or parasites (e.g., \textit{Plasmodium} species causing malaria, \textit{Giardia duodenalis}) must be excluded, particularly in travelers or immunocompromised hosts.
    \item \textbf{Systemic Inflammatory and Autoimmune Diseases}: Autoimmune conditions such as systemic lupus erythematosus (SLE), adult-onset Still's disease, and various vasculitides can present with prolonged fever, joint pain, fatigue, and rashes, closely mimicking acute or subacute viral syndromes.
    \textbf{Drug-Induced Hypersensitivity}: Adverse drug reactions, including drug-induced fever and cutaneous reactions (such as Stevens-Johnson syndrome or Drug Reaction with Eosinophilia and Systemic Symptoms [DRESS] syndrome), can mimic viral exanthems and systemic viral prodromes.
    \item \textbf{Hematological Malignancies}: Conditions such as acute leukemia or lymphoma frequently present with constitutional symptoms (fever, night sweats, weight loss), lymphadenopathy, and splenomegaly, which can be misdiagnosed as infectious mononucleosis (caused by Epstein-Barr virus) or other viral infections.
    \item \textbf{Non-Infectious Organ Dysfunction}: Acute non-infectious conditions, such as pulmonary embolism, myocardial infarction, or chemical pneumonitis, can present with acute respiratory distress, mimicking severe viral pneumonia or viral myocarditis.
\end{itemize}

\section*{When to Seek Medical Care}
While many viral infections are self-limiting and can be managed safely at home, certain clinical signs indicate severe disease or rapid deterioration, requiring prompt medical evaluation or emergency intervention.

Individuals should contact a healthcare professional or seek medical attention if they experience:
\begin{itemize}
    \item A high fever that remains unresponsive to over-the-counter antipyretic medications after 48 to 72 hours.
    \item Persistent vomiting or diarrhea that prevents the retention of oral fluids, leading to signs of moderate dehydration (e.g., severe thirst, dark urine, dry mouth, dizziness upon standing).
    \item The development of a progressive skin rash, particularly if accompanied by systemic symptoms.
    \item Symptoms that initially improve but then suddenly worsen, which may indicate a secondary bacterial infection.
\end{itemize}

Immediate emergency medical care must be sought if any of the following red flag symptoms occur:
\begin{itemize}
    \item Severe difficulty breathing, shortness of breath, persistent chest pain, or pressure.
    \item Confusion, severe lethargy, difficulty waking up, or altered mental status.
    \item A stiff neck, severe headache, and extreme sensitivity to light (photophobia).
    \item Blue, gray, or pale skin, lips, or nail beds (cyanosis), indicating inadequate oxygenation.
    \item A rash of dark red or purple spots that does not fade when pressed under a glass (a non-blanching purpuric rash).
    \item Seizures or sudden focal weakness.
\end{itemize}

In life-threatening situations, contact emergency services immediately by dialing:
\begin{itemize}
    \item \textbf{local emergency services} in many countries.
    \item \textbf{911} in the United States and Canada.
    \item \textbf{999} in the United Kingdom.
    \item \textbf{112} as the standard emergency number across the European Union.
\end{itemize}

\section*{Diagnosis and Investigations}
The diagnosis of viral infections relies on a combination of clinical evaluation and diagnostic testing. Identifying the specific causative virus guides clinical decision-making, epidemiological tracking, and the implementation of infection control measures.
\begin{itemize}
    \item \textbf{Clinical Assessment}: A comprehensive history, including the onset and progression of symptoms, travel history, immunization status, occupational exposures, and contact with sick individuals, is the cornerstone of diagnosis. This is paired with a physical examination to evaluate vital signs, respiratory effort, hydration status, lymph node enlargement, and mucocutaneous lesions.
    \item \textbf{Molecular Assays (PCR)}: Polymerase Chain Reaction (PCR) and other nucleic acid amplification tests (NAATs) are the gold standard for diagnosing many viral infections. These tests amplify specific viral DNA or RNA sequences, allowing for rapid and highly sensitive detection. PCR is widely used for respiratory panels, cerebrospinal fluid analysis, and blood-borne virus monitoring.
    \item \textbf{Serological Testing}: Serology involves detecting host antibodies (IgM and IgG) directed against viral antigens in serum or plasma. The presence of IgM antibodies typically indicates an acute, recent infection, while IgG antibodies suggest past exposure, recovery, or vaccination. Serology is essential for diagnosing hepatitis, HIV, and tracking herd immunity.
    \item \textbf{Antigen Detection}: Rapid antigen tests detect viral proteins directly from clinical specimens. While generally less sensitive than molecular tests, they provide rapid results (often within 15 minutes) and are highly useful for point-of-care screening of influenza, RSV, and COVID-19.
    \item \textbf{Hematological and Biochemical Profiles}: Routine blood tests provide supportive diagnostic clues. A Full Blood Count (FBC) may show leukopenia (low white blood cell count) or lymphocytosis (elevated lymphocytes), which are characteristic of viral infections. Inflammatory markers, such as C-Reactive Protein (CRP) and Procalcitonin (PCT), are typically normal or mildly elevated, helping to differentiate viral from bacterial infections.
    \item \textbf{Diagnostic Imaging}: For respiratory viral infections, a chest X-ray or high-resolution computed tomography (HRCT) scan may demonstrate ground-glass opacities, consolidation, or interstitial infiltrates. In cases of suspected viral encephalitis, magnetic resonance imaging (MRI) of the brain can reveal characteristic inflammatory changes in the temporal lobes or other brain regions.
\end{itemize}

\section*{Management}
The management of viral infections is primarily supportive, aimed at relieving symptoms and maintaining physiological stability while the host's immune system clears the pathogen. However, specific pharmacological interventions are available for certain viral diseases.
\begin{itemize}
    \item \textbf{Supportive Care}: The foundation of viral management includes absolute rest to conserve energy, and aggressive hydration to replace fluid losses from fever, sweating, vomiting, or diarrhea. Oral rehydration solutions (ORS) are preferred for gastrointestinal infections to maintain electrolyte balance.
    \item \textbf{Symptomatic Pharmacotherapy}: Over-the-counter medications are utilized to control symptoms. Analgesics and antipyretics, such as paracetamol (acetaminophen) or ibuprofen, help manage fever, myalgia, and headache. It is critical to avoid aspirin in children and adolescents presenting with viral symptoms due to the risk of Reye's syndrome—a rare but potentially fatal condition causing acute encephalopathy and liver dysfunction.
    \item \textbf{Antiviral Agents}: Targeted antiviral therapies inhibit specific stages of the viral replication cycle. Examples include neuraminidase inhibitors (e.g., oseltamivir) for influenza, nucleoside analogs (e.g., acyclovir, valacyclovir) for herpesviruses, and protease inhibitors or polymerase inhibitors (e.g., nirmatrelvir/ritonavir) for COVID-19. For chronic infections like HIV, lifelong antiretroviral therapy (ART) is required to suppress viral replication and preserve immune function.
    \item \textbf{Immunomodulatory Therapies}: In cases of severe viral infections characterized by hyperinflammation, immunomodulatory agents may be indicated. Corticosteroids (e.g., dexamethasone) have been shown to reduce mortality in patients with severe COVID-19 requiring oxygen. Intravenous immunoglobulin (IVIG) may be used to provide passive immunity or modulate immune-mediated complications.
    \item \textbf{Critical Care Support}: For severe cases resulting in respiratory failure (e.g., ARDS) or cardiovascular collapse, intensive care unit (ICU) admission is required. Management may involve supplemental oxygen therapy, non-invasive positive pressure ventilation, mechanical ventilation, or extracorporeal membrane oxygenation (ECMO).
\end{itemize}

\section*{Complications}
Viral infections can lead to serious secondary complications, either during the acute phase of the illness or as long-term sequelae. These complications can affect multiple organ systems and significantly impact patient outcomes.
\begin{itemize}
    \item \textbf{Secondary Bacterial Superinfections}: Viral damage to mucosal barriers (such as the respiratory epithelium) compromises host defenses, making the patient susceptible to secondary bacterial infections. Examples include bacterial pneumonia (commonly caused by \textit{Streptococcus pneumoniae} or \textit{Staphylococcus aureus}) following an influenza infection, and secondary bacterial skin infections following varicella or monkeypox lesions.
    \item \textbf{Chronic Infection, Latency, and Oncogenesis}: Certain viruses establish persistent infections. Hepatitis B and C can cause chronic active hepatitis, leading to liver cirrhosis and hepatocellular carcinoma. High-risk strains of Human Papillomavirus (HPV) are oncogenic, directly linked to cervical, anogenital, and oropharyngeal cancers. Latent viruses, such as varicella-zoster, can reactivate decades later to cause herpes zoster (shingles) and postherpetic neuralgia.
    \item \textbf{Post-Viral Inflammatory and Autoimmune Syndromes}: Viral antigens can trigger aberrant immune responses that target host tissues. This includes Guillain-Barr\'{e} syndrome (an acute demyelinating polyneuropathy), reactive arthritis, and post-viral fatigue syndrome. Long COVID (post-acute sequelae of COVID-19) is a multi-system syndrome characterized by persistent fatigue, dyspnoea, and cognitive dysfunction.
    \item \textbf{Systemic and Organ-Specific Injury}: Severe viral infections can lead to acute myocarditis (inflammation of the heart muscle), pericarditis, acute kidney injury (AKI), hepatotoxicity, and viral encephalopathy. Severe systemic involvement can culminate in disseminated intravascular coagulation (DIC) and septic shock.
\end{itemize}

\section*{Prognosis and Outcomes}
The prognosis of a viral infection is highly variable, depending on a combination of host factors, viral virulence, and the availability of medical interventions. The vast majority of acute viral infections (such as the common cold, mild gastroenteritis, and uncomplicated influenza) are self-limiting, with patients achieving full recovery within 7 to 14 days without long-term sequelae.

For chronic or severe viral infections, the prognosis has dramatically improved due to modern medical advances. For example:
\begin{itemize}
    \item The introduction of highly active antiretroviral therapy (HAART) has transformed HIV/AIDS from a uniformly fatal disease into a manageable, chronic health condition with a near-normal life expectancy for patients adhering to treatment.
    \item Direct-acting antiviral (DAA) therapies for chronic hepatitis C now achieve cure rates exceeding 95\%, significantly reducing the risk of progressive liver failure and liver cancer.
\end{itemize}

Conversely, prognosis remains guarded for severe infections caused by highly pathogenic viruses, such as Ebola virus, rabies (which is virtually 100\% fatal once clinical symptoms appear), and severe cases of viral encephalitis or hemorrhagic fevers. Key determinants of a poor prognosis include advanced age, the presence of pre-existing cardiovascular or pulmonary comorbidities, severe immunosuppression, delayed presentation to healthcare facilities, and lack of access to supportive care or specific antiviral therapies. Long-term outcomes for survivors of severe acute infections may be complicated by prolonged physical rehabilitation requirements, pulmonary fibrosis, or cognitive deficits.

\section*{Prevention and Self-Care}
Preventing the transmission of viral pathogens and managing mild symptoms effectively at home are key strategies for reducing the individual and societal burden of viral diseases.
\begin{itemize}
    \item \textbf{Immunization (Vaccination)}: Vaccination is the most effective public health intervention for preventing viral infections and their complications. Vaccines work by stimulating the immune system to produce antibodies and memory cells without causing the disease itself. Examples include the measles, mumps, and rubella (MMR) vaccine, seasonal influenza vaccines, hepatitis B vaccines, human papillomavirus (HPV) vaccines, and COVID-19 vaccines. Achieving high vaccination rates is critical for establishing herd immunity, which protects individuals who cannot be vaccinated.
    \item \textbf{Hand and Personal Hygiene}: Regular hand hygiene is a simple yet powerful preventive measure. Hands should be washed thoroughly with soap and water for at least 20 seconds, especially after coughing, sneezing, using the restroom, or before eating. Alcohol-based hand sanitizers containing at least 60\% alcohol can be used when soap and water are unavailable. Good respiratory hygiene involves covering the mouth and nose with a tissue or the crook of the elbow when coughing or sneezing.
    \item \textbf{Environmental Cleaning and Sanitation}: Disinfecting frequently touched surfaces (such as doorknobs, countertops, and mobile phones) with registered disinfectants kills viral particles on surfaces (fomites). Ensuring access to clean drinking water, proper food handling, and robust sewage treatment facilities prevents the transmission of waterborne and foodborne viruses (e.g., norovirus, hepatitis A).
    \item \textbf{Physical Distancing and Protective Barriers}: During outbreaks or in high-transmission settings, wearing face masks (such as surgical masks or N95/P2 respirators) reduces the inhalation and exhalation of infectious respiratory droplets. Maintaining physical distance from individuals showing symptoms of illness further minimizes exposure.
    \item \textbf{Safe Sexual and Injection Practices}: The use of barrier methods (condoms) during sexual intercourse reduces the transmission of sexually transmitted viruses, including HIV, herpes simplex virus, and human papillomavirus. Needle exchange programs and the use of sterile syringes prevent the transmission of blood-borne viruses (hepatitis B, hepatitis C, HIV) among people who inject drugs.
    \item \textbf{Vector Avoidance}: To prevent vector-borne viral infections, individuals should apply insect repellents containing DEET, picaridin, or oil of lemon eucalyptus, wear long-sleeved clothing, sleep under insecticide-treated bed nets, and eliminate standing water around living spaces to prevent vector breeding.
    \item \textbf{Home Care and Self-Monitoring}: When managing a mild viral infection at home, individuals should prioritize rest, maintain hydration by drinking water, herbal teas, or clear broths, and self-isolate to prevent infecting household members. Monitoring symptoms (such as body temperature, respiratory rate, and oxygen saturation using a pulse oximeter) is crucial for identifying early signs of clinical deterioration.
\end{itemize}

\section*{Sources and Further Reading}
\begin{itemize}
    \item World Health Organization (WHO): Disease Outbreak News and viral infection factsheets. URL: \url{https://www.who.int}
    \item Centers for Disease Control and Prevention (CDC): Comprehensive resources on viral diseases, transmission, and prevention. URL: \url{https://www.cdc.gov}
    \item National Institutes of Health (NIH) - National Institute of Allergy and Infectious Diseases (NIAID): Detailed scientific summaries on viral pathogenesis and research. URL: \url{https://www.niaid.nih.gov}
    \item NHS Choices: Clinical overview of common viral infections, management, and patient guidance. URL: \url{https://www.nhs.uk}
    \item Australian Government Department of Health and Aged Care: National guidelines on immunization and communicable disease control. URL: \url{https://www.health.gov.au}
    \item Mayo Clinic: Patient-oriented resources on viral symptoms, diagnosis, and treatment. URL: \url{https://www.mayoclinic.org}
\end{itemize}